\section{Вступ}\label{sec:vstup}

\subsection{Призначення та цільова аудиторія}\label{sec:vstup-meta}
Цей документ являє собою внутрішню специфікацію вимог до програмного забезпечення (Software Requirements Specification, SRS) для інформаційної системи автоматизації бронювання квитків у кінотеатрі (Cinema Booking System)~\sked{K6}. Специфікація орієнтована на формат стандарту IEEE 830 (із врахуванням рекомендацій ISO/IEC/IEEE 29148) та призначена для фіксації погоджених вимог і обмежень проєкту між автором та автономним AI-агентом розробки Google Antigravity~\sked{K6,K10}.

Документ є єдиним джерелом правди для подальшої практичної реалізації системи за методологією Test-Driven Development (TDD) та верифікації її коректності за допомогою автоматизованих тестів~\sked{K10}.

\subsection{Межі системи та виключення версії 1}\label{sec:vstup-mezhi}
Система призначена для автоматизації процесів вибору сеансів та резервування місць у залах кінотеатру для двох ключових ролей: Користувача та Адміністратора~\sked{K1}. 

Для забезпечення високої надійності та швидкого постачання першої робочої версії (v1) низка функціональних можливостей свідомо винесена за межі системи:
\begin{itemize}
    \item \textbf{Гостьове бронювання та стороння автентифікація:} створення та скасування бронювань вимагають обов'язкової реєстрації у системі; гостьове бронювання та інтеграція із зовнішніми провайдерами автентифікації (OAuth, SSO, Google тощо) у першій версії не передбачені~\sked{K25}.
    \item \textbf{Оплата та фінансові транзакції:} система реалізує виключно безоплатне резервування місць; онлайн-оплата, інтеграція з платіжними сервісами (шлюзами), відстеження статусів оплати та процедури повернення коштів винесені за межі інформаційної системи~\sked{K33}.
    \item \textbf{Вартість квитків:} ціни квитків у першій версії не зберігаються, не обробляються базою даних та не відображаються у вебінтерфейсі користувача чи адміністратора~\sked{K34}.
    \item \textbf{Оплата квитків офлайн:} фактичне придбання квитків та їх оплата здійснюються користувачами безпосередньо в касах або терміналах кінотеатру поза межами цієї програми~\sked{K35}.
\end{itemize}

\subsection{Глосарій, терміни та скорочення}\label{sec:vstup-hlosarii}
У цій специфікації використовуються терміни та скорочення, наведені в \tabref{tab:glossary}~\sked{K6,K10}.

\begin{table}[htbp]
\caption{Глосарій ключових термінів та скорочень}\label{tab:glossary}
\centering
\begin{tabularx}{\textwidth}{@{}l l L@{}}
\toprule
\textbf{Термін/Скорочення} & \textbf{Повна назва} & \textbf{Визначення в контексті проєкту} \\
\midrule
SRS & Software Requirements Specification & Специфікація вимог до ПЗ відповідно до формату IEEE 830. \\
TDD & Test-Driven Development & Інженерна методологія розробки через попереднє написання автоматизованих тестів. \\
BDD & Behavior-Driven Development & Опис поведінки системи мовою бізнес-сценаріїв (Given-When-Then / Gherkin). \\
RBAC & Role-Based Access Control & Модель розмежування прав доступу на основі ролей облікового запису (Користувач, Адміністратор). \\
Фільм & Movie / Film & Кінокартина або художній твір, що має назву, тривалість, опис і демонструється в кінотеатрі. \\
Зал & Cinema Hall & Приміщення кінотеатру з визначеною назвою, місткістю та схемою розташування місць. \\
Сеанс & Showtime & Запланований показ фільму в залі з точним часом початку ($T_{сеансу}$). \\
Місце & Seat & Окрема одиниця розміщення глядача в залі (ряд і номер місця). \\
Бронювання & Booking / Reservation & Фіксація прав зареєстрованого користувача на перелік місць (від 1 до 6 включно) на обраний сеанс. \\
Атомарність & Atomicity & Властивість транзакції: створення бронювання виконується неподільно («все або нічого»). \\
Hold / Таймаут & Temporary Lock & Механізм тимчасового утримання місць на етапі вибору; у версії 1 не застосовується. \\
\bottomrule
\end{tabularx}
\end{table}

\subsection{Нормативні посилання та огляд структури}\label{sec:vstup-struktura}
Специфікація спирається на базові засади стандарту IEEE 830-1998 «IEEE Recommended Practice for Software Requirements Specifications» та проєктні домовленості, зафіксовані в репозиторії~\sked{K6}. 

Документ структуровано наступним чином: Розділ~\ref{sec:opis} наводить загальний архітектурний опис та обмеження; Розділ~\ref{sec:interfeisy} визначає контракт зовнішніх інтерфейсів на рівні операцій; Розділ~\ref{sec:funktsional} специфікує функціональні вимоги (REQ-F-xxx); Розділ~\ref{sec:tdd} встановлює обов'язковий регламент TDD і BDD-сценарії; Розділ~\ref{sec:loguvannia} регламентує правила аудиту та логування; Розділ~\ref{sec:nefunktsional} охоплює вимоги до якості та атомарності; Розділ~\ref{sec:vidkladeni} фіксує відкладені рішення щодо стеку; Розділ~\ref{sec:vysnovky} підсумовує готовність до розробки, а Додаток~\ref{sec:dodatok-matrytsia} містить матрицю простежуваності.
